\documentclass[aspectratio=169]{beamer}
\usetheme{Madrid}
\usecolortheme{dolphin}
\setbeamertemplate{navigation symbols}{}
\setbeamertemplate{footline}[frame number]
\setbeamersize{text margin left=7mm,text margin right=7mm}
\usepackage[T1]{fontenc}
\usepackage{lmodern}
\usepackage{booktabs,tabularx,array}
\usepackage{xcolor}
\hypersetup{hidelinks,pdftitle={Allocating Bits Between Coordinate Pairs}}
\newcommand{\good}[1]{\textcolor{green!45!black}{\textbf{#1}}}
\newcommand{\bad}[1]{\textcolor{red!75!black}{\textbf{#1}}}
\title{Allocating Bits Between Coordinate Pairs}
\subtitle{A real signal before SAQ rotation, but no usable opportunity after it}
\author{}
\date{2026-09-04}

\begin{document}
\begin{frame}\titlepage\end{frame}

\begin{frame}{1. The Answer First}
\begin{alertblock}{Final decision}
Stop the current pair-level allocation direction. Do not build its query
consumer.
\end{alertblock}
\begin{itemize}
  \item Before SAQ's internal rotation, moving bits between coordinate pairs
  looks valuable: about 28--31\% on SIFT and 6.3--6.4\% on GIST residuals.
  \item After the rotation used inside each SAQ segment, every pair chooses
  the bit width it already had.
  \item At exactly matched payload and factor bytes, the final measured
  direction-error improvement is \bad{0\%}.
\end{itemize}
\begin{block}{Main lesson}
The earlier positive result was not numerically wrong. It measured structure
that a later SAQ step removes before quantization and search.
\end{block}
\end{frame}

\begin{frame}{2. Correction: What Was the Intended Question?}
\begin{columns}[T]
\column{0.47\textwidth}
\begin{block}{Earlier interpretation}
For one pair $(D_1,D_2)$, decide how many states belong to $D_1$ and how many
belong to $D_2$.
\end{block}
\column{0.47\textwidth}
\begin{block}{The intended question}
Treat $(D_1,D_2)$ as one unit and $(D_3,D_4)$ as another. Decide how much of
the total storage each \emph{pair} receives.
\end{block}
\end{columns}
\begin{alertblock}{Why this matters}
This is allocation \textbf{between coordinate pairs}, not allocation between
the two coordinates inside one pair. The September branch was created to test
this corrected question directly.
\end{alertblock}
\end{frame}

\begin{frame}{3. Running Example: Give Bits to the Harder Pair}
\small
Assume two coordinate pairs share a total budget of 16 bits.
\begin{center}
\begin{tabular}{lccc}
\toprule
Pair & Variation in the data & Equal allocation & Flexible allocation\\
\midrule
$(D_1,D_2)$ & large & 8 bits & 10 bits\\
$(D_3,D_4)$ & small & 8 bits & 6 bits\\
\midrule
Total & & 16 bits & 16 bits\\
\bottomrule
\end{tabular}
\end{center}
\begin{block}{Intuition}
An extra bit creates more reconstruction levels. It should go to the pair
where those extra levels reduce more error. The easy pair gives up two bits;
the hard pair receives them; total storage does not change.
\end{block}
\begin{alertblock}{Testable hypothesis}
If different pairs have persistently different returns from an extra bit,
flexible allocation should lower held-out error by at least 5\%.
\end{alertblock}
\end{frame}

\begin{frame}{4. Where Does This Happen in the SAQ Pipeline?}
\small
\begin{center}
\begin{tabular}{c}
original vector\\[-1mm]
$\downarrow$\\[-1mm]
principal-component coordinates\\[-1mm]
$\downarrow$\\[-1mm]
subtract the coarse-region center to obtain a residual\\[-1mm]
$\downarrow$\\[-1mm]
split the residual into SAQ segments\\[-1mm]
$\downarrow$\\[-1mm]
\alert{rotate the coordinates inside each segment}\\[-1mm]
$\downarrow$\\[-1mm]
quantize and store bits
\end{tabular}
\end{center}
\begin{block}{Two places to measure}
The first diagnostic measured pairs before the segment rotation. The decisive
check measured the coordinates that the SAQ quantizer actually receives,
after that rotation.
\end{block}
\scriptsize A residual is simply the vector left after subtracting the center
of its coarse database region.
\end{frame}

\begin{frame}[fragile]{5. First Diagnostic: How We Allocated the Budget}
\footnotesize
\begin{verbatim}
Use 16,384 fixed base vectors; split them into halves A and B.

For each representation stage and each two-coordinate group:
  fit the local quantizer on A
  measure error for candidate group budgets from 6 to 10 bits

Use exact dynamic programming:
  choose one budget for every group
  keep the total at exactly 512 bits
  minimize total fitting error

Evaluate the frozen allocation on B; then swap A and B.
\end{verbatim}
\begin{alertblock}{No query tuning}
Only learn/base vectors were used. Benchmark queries, exact nearest-neighbour
answers, Recall, and query speed were not read.
\end{alertblock}
\end{frame}

\begin{frame}{6. The Signal Survived PCA and Coarse Residualization}
\small
Held-out reconstruction-error reduction from allocating bits between fixed
adjacent pairs:
\begin{center}
\begin{tabular}{lccc}
\toprule
Dataset & principal-component view & 1,024 coarse regions & 4,096 coarse regions\\
\midrule
SIFT1M & 40.59\% & 22.55\% & 20.11\%\\
GIST1M & 30.82\% & 13.93\% & 11.73\%\\
\bottomrule
\end{tabular}
\end{center}
\begin{block}{What this established}
Unequal demand between two-coordinate groups is not merely an artifact of the
original raw coordinates. It remains visible after the main representation
steps used to prepare SAQ residuals.
\end{block}
\begin{alertblock}{What this did not establish}
The comparison was against equal 8-bit group budgets, not yet against the
complete SAQ representation.
\end{alertblock}
\end{frame}

\begin{frame}{7. Correlation Was Not the Main Explanation}
\small
For each residual group, we compared the value of one extra bit with two
simple statistics.
\begin{center}
\begin{tabularx}{0.88\textwidth}{>{\bfseries}p{0.35\textwidth}X}
Total variation of the pair & correlation with extra-bit value was
0.966--0.995\\
Absolute correlation between the coordinates & only 0.041--0.089 on GIST and
0.246--0.306 on SIFT\\
\end{tabularx}
\end{center}
\begin{alertblock}{Interpretation}
The useful rule was essentially ``give more bits to groups containing more
energy.'' It was not ``pair the most strongly correlated coordinates.''
\end{alertblock}
This is important because ordinary rate--distortion allocation already uses
the first principle.
\end{frame}

\begin{frame}{8. Closest Work Already Covers the General Principle}
\footnotesize
\begin{tabularx}{\textwidth}{>{\bfseries}p{0.25\textwidth}X}
Transform Coding (2010) & principal-component transform, trained scalar
quantizers, and unequal bit allocation under one total budget\\
Optimized Transform Coding (2014) & component-specific error estimates for
choosing the bit allocation\\
OPQ (2013) & optimized rotation and decomposition for product quantization\\
Adaptive Bit Allocation PQ (2016) & unequal codebook sizes between product-
quantizer subspaces\\
Distribution-Sensitive PQ (2018) & moves bits between subspaces according to
their distributions\\
SAQ (2026) & dynamic programming for segment boundaries and one bit width per
segment under a storage budget\\
\end{tabularx}
\begin{alertblock}{Novelty decision}
Local 2D rotation + fitted error curves + exact budget DP is a direct
composition of known ideas. Exact DP improves correctness; it does not create
a new quantization model.
\end{alertblock}
\end{frame}

\begin{frame}{9. The Narrow SAQ-Specific Question}
\begin{block}{What might still have been new}
Inside one existing SAQ segment, let different coordinate pairs store
different numbers of lower bits, while preserving SAQ's shared scale and
two-stage scan.
\end{block}
The test had to preserve all of the following:
\begin{itemize}
  \item the same segment boundaries and number of segment factors;
  \item exactly the same total payload and factor bytes;
  \item one common scale for the whole vector segment;
  \item one fast bit for every represented coordinate; and
  \item no per-vector plan identifier or per-pair factor.
\end{itemize}
\begin{alertblock}{Why test before implementation?}
A ragged-bit query consumer would add packing, decoding, and branching work.
We first asked whether enough base-only quality opportunity remained to justify
that engineering.
\end{alertblock}
\end{frame}

\begin{frame}{10. The Exact Storage-Matched Baselines}
\small
\begin{center}
\begin{tabularx}{\textwidth}{lXr}
\toprule
Dataset & Existing SAQ plan & Total counted storage\\
\midrule
SIFT1M & all 128 dimensions use 4 bits each & 72 bytes/vector\\
GIST1M & successive groups of 64, 192, 320, 256, and 128 dimensions use
11, 6, 4, 2, and 0 bits each & 488 bytes/vector\\
\bottomrule
\end{tabularx}
\end{center}
\begin{block}{Why factors are counted}
Every positive SAQ segment stores two additional floating-point values per
database vector. Allowing more segments would secretly spend more bytes, so
we kept the existing segments and factor count unchanged.
\end{block}
Only widths \emph{inside} each positive segment could move between adjacent
pairs.
\end{frame}

\begin{frame}{11. Running Example: Rotation Removes the Imbalance}
\small
Consider four independent coordinates grouped into two pairs:
\[
\underbrace{(D_1,D_2)}_{\text{variances }9,9},\qquad
\underbrace{(D_3,D_4)}_{\text{variances }1,1}.
\]
Before rotation, the first pair contains 90\% of the total variation, so it
clearly deserves more bits.
\medskip

A balanced orthogonal rotation mixes all four coordinates without losing
information. In this simple example, it can spread the same total variation
as
\[
(5,5,5,5).
\]
\begin{alertblock}{After rotation}
Both new pairs now have equal variation. Moving a bit from one pair to the
other only makes one side worse; the equal-width allocation is optimal again.
\end{alertblock}
SAQ performs this kind of mixing separately inside every segment.
\end{frame}

\begin{frame}[fragile]{12. Decisive Opportunity Decomposition}
\footnotesize
\begin{verbatim}
Recover the current SAQ plan from fitting-half variances.

For every positive segment:
  keep its dimensions, payload bits, and two stored factors
  apply one fixed and recorded orthogonal rotation
  fit actual uniform-lattice error curves for every adjacent pair
  reallocate pair widths with exact dynamic programming

For held-out vectors:
  use the segment-wide maximum and shared scale
  run the same six rounds of coordinate adjustment
  compare direction-related error with the uniform segment baseline

Swap the fitting and evaluation halves and repeat.
\end{verbatim}
\begin{block}{Frozen stopping rule}
Stop if the post-rotation gain is below 5\% on either natural dataset.
\end{block}
\end{frame}

\begin{frame}{13. Before Rotation: The Opportunity Still Looks Large}
\small
Held-out improvement predicted by SAQ's own variance-based planning formula:
\begin{center}
\begin{tabular}{lcc}
\toprule
Dataset and residual view & A fitted, B tested & B fitted, A tested\\
\midrule
SIFT, 1,024 regions & 30.74\% & 30.83\%\\
SIFT, 4,096 regions & 28.20\% & 28.39\%\\
GIST, 1,024 regions & 6.45\% & 6.39\%\\
GIST, 4,096 regions & 6.39\% & 6.31\%\\
\bottomrule
\end{tabular}
\end{center}
\begin{alertblock}{This is an upper-level diagnostic, not the answer}
These coordinates have not yet passed through the rotation immediately before
SAQ quantization. Using this table alone would repeat the earlier mistake.
\end{alertblock}
\end{frame}

\begin{frame}{14. After Rotation: Every Pair Returns to Baseline}
\small
\begin{center}
\begin{tabular}{lccc}
\toprule
Dataset & residual settings & pairs changing width & final error gain\\
\midrule
SIFT1M & 2 region counts $\times$ 2 folds & \bad{0} & \bad{0\%}\\
GIST1M & 2 region counts $\times$ 2 folds & \bad{0} & \bad{0\%}\\
\bottomrule
\end{tabular}
\end{center}
\begin{itemize}
  \item Every SIFT pair selects the existing 4-bit width.
  \item Every positive GIST pair selects its segment's existing 11-, 6-, 4-,
  or 2-bit width.
  \item The variance formula and the fitted uniform-lattice curves agree.
  \item With identical selected widths, the adjusted database codes and their
  measured error are identical to the baseline.
\end{itemize}
\begin{alertblock}{Gate result}
The required 5\% opportunity is missed by the full 5 percentage points.
\end{alertblock}
\end{frame}

\begin{frame}{15. Why the Two Positive Tables Were Not Contradictory}
\begin{columns}[T]
\column{0.47\textwidth}
\begin{block}{Before segment rotation}
PCA orders coordinates from high to low importance. Adjacent pairs inherit
large differences, so bit allocation helps.
\end{block}
\column{0.47\textwidth}
\begin{block}{After segment rotation}
SAQ deliberately mixes coordinates inside each segment. Pair-level demand is
equalized, so the same allocation no longer helps.
\end{block}
\end{columns}
\begin{alertblock}{Research interpretation}
The signal exists at one stage of the pipeline but is not available to the
unchanged SAQ representation and consumer. Representation-stage attribution
prevented us from building the wrong system.
\end{alertblock}
\end{frame}

\begin{frame}{16. What Is Closed, and What Is Not?}
\begin{columns}[T]
\column{0.47\textwidth}
\begin{block}{Closed by this result}
\begin{itemize}
  \item keep the current SAQ segment rotation;
  \item reallocate bits between adjacent pairs afterward;
  \item build a ragged mixed-width consumer for that allocation.
\end{itemize}
\end{block}
\column{0.47\textwidth}
\begin{block}{Not proved impossible}
\begin{itemize}
  \item every possible orthogonal rotation;
  \item jointly redesigning transform and allocation;
  \item unrelated SAQ-specific limitations.
\end{itemize}
\end{block}
\end{columns}
\begin{alertblock}{Important novelty boundary}
A joint transform/allocation method would be a new research question under
strong Transform Coding, OPQ, BAPQ, and DSPQ prior-work pressure. It is not a
positive continuation supported by this experiment.
\end{alertblock}
\end{frame}

\begin{frame}{17. Evidence Quality and Cost}
\small
\begin{itemize}
  \item Official TexMex SIFT1M and GIST1M learn/base files were hash verified.
  \item No benchmark query, exact-answer file, Recall result, QPS result, or
  old ANN index result was read.
  \item The two data halves were used in both fit/evaluation directions.
  \item Ten focused tests pass; generalized mixed-width adjustment matches an
  independent scalar implementation on a tiny fixture.
  \item Two accepted executions produced byte-identical metadata and complete
  scientific summaries.
  \item Each accepted run used about 94 CPU-seconds and less than 716 MB peak
  memory.
\end{itemize}
\begin{block}{Performance status}
This was an early base-only falsification test, not Recall or query-speed
evidence. Because the quality opportunity is zero, production optimization is
not justified.
\end{block}
\end{frame}

\begin{frame}{18. Final Takeaway}
\begin{alertblock}{Decision}
Stop the current post-rotation pair-allocation direction.
\end{alertblock}
\begin{enumerate}
  \item The corrected pair-level question was tested directly.
  \item Unequal group demand is real before the final SAQ segment rotation.
  \item Variance, not within-pair correlation, explains most of the signal.
  \item Existing work already covers the general allocation principle.
  \item Under the unchanged SAQ representation and exact byte budget, the
  remaining usable opportunity is 0\%.
\end{enumerate}
\begin{block}{Next research step}
Close this branch. Start a new direction only from a limitation specific to
SAQ's shared factors, adjustment rule, or two-stage query consumer, rather
than trying to rescue this allocation result.
\end{block}
\end{frame}

\begin{frame}{References}
\scriptsize
\begin{thebibliography}{9}
\bibitem{brandt} Brandt. \emph{Transform Coding for Fast Approximate Nearest
Neighbor Search in High Dimensions}. CVPR, 2010.
\bibitem{opq} Ge et al. \emph{Optimized Product Quantization for Approximate
Nearest Neighbor Search}. CVPR, 2013.
\bibitem{otc} Park et al. \emph{Optimized Transform Coding for Approximate
KNN Search}. BMVC, 2014.
\bibitem{bapq} Guo et al. \emph{Adaptive Bit Allocation Product
Quantization}. Neurocomputing, 2016.
\bibitem{dspq} Li et al. \emph{Distribution Sensitive Product Quantization}.
IEEE TCSVT, 2018.
\bibitem{quicker} Andr\'e et al. \emph{Quicker ADC: Unlocking the Hidden
Potential of Product Quantization with SIMD}. IEEE TPAMI, 2021.
\bibitem{saq} Li et al. \emph{SAQ: Pushing the Limits of Vector Quantization
through Code Adjustment and Dimension Segmentation}. ACM SIGMOD, 2026;
arXiv:2509.12086.
\end{thebibliography}
\vfill
\tiny Detailed project evidence: representation-stage attribution,
closest-primary-work review, and matched-byte opportunity decomposition dated
2026-09-03 in \texttt{docs/research/}.
\end{frame}

\end{document}
